% Nota de dissertacao_pb. GERADA por flim_al/tabelas.py.
\footnotesize \texttt{random} é o FLIM puro: sorteia as imagens, como o artigo faz na primeira rodada. Todos os braços recebem o MESMO número de imagens (K) e são medidos no mesmo conjunto de teste. \textbf{\texttt{regiao\_confusa} é diferente dos outros três}: ele usa exatamente as imagens que o sorteio usaria naquela semente e muda só ONDE o traço cai, com o mesmo número de pixels --- então o $\Delta$ dele mede posição de anotação, e o dos outros mede escolha de imagem. \textbf{Acurácia é \texttt{1 - MAE}}, a fração de pixels certos; ela passa de 0,97 em quase tudo porque o objeto ocupa uma fração mínima da imagem, e acertar o fundo já garante isso --- quem separa os braços é F$\beta$ e IoU. $\Delta$ e p são pareados por semente contra o FLIM puro, no mesmo dataset e no mesmo K; com poucas sementes o teste tem pouco poder, e p alto significa \textbf{não decidido}, não 'igual'. Decoder FLIM\_pb.
\normalsize
